\Needspace{20\baselineskip}\SecondArticle{16}
\textsc{Theorem.} \textit{Let $\varphi(u,x)$}\HistoricalNote{12}\textit{ denote a product of any number of factors in which the indeterminates $u,x$ enter only linearly, that is, factors of the form}
\[
\begin{gathered}
\alpha+\beta u+\gamma x\\
\alpha'+\beta'u+\gamma'x\\
\alpha''+\beta''u+\gamma''x\\
\text{etc.}
\end{gathered}
\]
\textit{and let $w$ be another indeterminate. Then the function}
\[
\varphi(u+w\, .\,\tfrac{\mathrm d\varphi(u,x)}{\mathrm dx},\ x-w\, .\,\tfrac{\mathrm d\varphi(u,x)}{\mathrm du})=\Omega
\]
\textit{is identically divisible by $\varphi(u,x)$.}



\SourceLeaf{059}
\textit{Proof.} Setting
\[
\begin{aligned}
\varphi(u,x)&=(\alpha+\beta u+\gamma x)Q\\
&=(\alpha'+\beta'u+\gamma'x)Q'\\
&=(\alpha''+\beta''u+\gamma''x)Q''\\
&\quad\text{etc.}
\end{aligned}
\]
we obtain for $Q,Q',Q''$ etc. polynomials in the indeterminates $u,\allowbreak x,\allowbreak \alpha,\allowbreak \beta,\allowbreak \gamma,\allowbreak \alpha',\allowbreak \beta',\allowbreak \gamma'$, $\alpha'',\beta'',\gamma''$ etc., and
\[
\begin{aligned}
\frac{\mathrm d\varphi(u,x)}{\mathrm dx}&=\gamma Q+(\alpha+\beta u+\gamma x)\, .\,\frac{\mathrm dQ}{\mathrm dx}\\
&=\gamma'Q'+(\alpha'+\beta'u+\gamma'x)\, .\,\frac{\mathrm dQ'}{\mathrm dx}\\
&=\gamma''Q''+(\alpha''+\beta''u+\gamma''x)\, .\,\frac{\mathrm dQ''}{\mathrm dx}\\
&\quad\text{etc.}
\end{aligned}
\]
\[
\begin{aligned}
\frac{\mathrm d\varphi(u,x)}{\mathrm du}&=\beta Q+(\alpha+\beta u+\gamma x)\, .\,\frac{\mathrm dQ}{\mathrm du}\\
&=\beta'Q'+(\alpha'+\beta'u+\gamma'x)\, .\,\frac{\mathrm dQ'}{\mathrm du}\\
&=\beta''Q''+(\alpha''+\beta''u+\gamma''x)\, .\,\frac{\mathrm dQ''}{\mathrm du}\\
&\quad\text{etc.}
\end{aligned}
\]
When these values are substituted in the factors whose product is $\Omega$, namely in
\[
\begin{gathered}
\alpha+\beta u+\gamma x+\beta w\, .\,\frac{\mathrm d\varphi(u,x)}{\mathrm dx}-\gamma w\, .\,\frac{\mathrm d\varphi(u,x)}{\mathrm du}\\
\alpha'+\beta'u+\gamma'x+\beta'w\, .\,\frac{\mathrm d\varphi(u,x)}{\mathrm dx}-\gamma'w\, .\,\frac{\mathrm d\varphi(u,x)}{\mathrm du}\\
\alpha''+\beta''u+\gamma''x+\beta''w\, .\,\frac{\mathrm d\varphi(u,x)}{\mathrm dx}-\gamma''w\, .\,\frac{\mathrm d\varphi(u,x)}{\mathrm du}\\
\text{etc., respectively}
\end{gathered}
\]
these take the following values:
\[
\begin{gathered}
(\alpha+\beta u+\gamma x)(1+\beta w\, .\,\tfrac{\mathrm dQ}{\mathrm dx}-\gamma w\, .\,\tfrac{\mathrm dQ}{\mathrm du})\\
(\alpha'+\beta'u+\gamma'x)(1+\beta'w\, .\,\tfrac{\mathrm dQ'}{\mathrm dx}-\gamma'w\, .\,\tfrac{\mathrm dQ'}{\mathrm du})\\
(\alpha''+\beta''u+\gamma''x)(1+\beta''w\, .\,\tfrac{\mathrm dQ''}{\mathrm dx}-\gamma''w\, .\,\tfrac{\mathrm dQ''}{\mathrm du})\\
\text{etc.}
\end{gathered}
\]
\SourceLeaf{060}
Thus $\Omega$ is the product of $\varphi(u,x)$ and the factors
\[
\begin{gathered}
1+\beta w\, .\,\frac{\mathrm dQ}{\mathrm dx}-\gamma w\, .\,\frac{\mathrm dQ}{\mathrm du}\\
1+\beta'w\, .\,\frac{\mathrm dQ'}{\mathrm dx}-\gamma'w\, .\,\frac{\mathrm dQ'}{\mathrm du}\\
1+\beta''w\, .\,\frac{\mathrm dQ''}{\mathrm dx}-\gamma''w\, .\,\frac{\mathrm dQ''}{\mathrm du}
\end{gathered}
\]
etc., that is, of $\varphi(u,x)$ and a polynomial in the indeterminates $u,x,w,\alpha,\beta,\gamma$, $\alpha',\beta',\gamma',\alpha'',\beta'',\gamma''$ etc. Q. E. D.

\SecondArticle{17}
The theorem of the preceding article clearly applies to the function $\zeta$, which we shall now represent by
\[
f(u,x,\lambda',\lambda'',\lambda'''\text{ etc.})
\]
so that
\[
f(u+w\, .\,\tfrac{\mathrm d\zeta}{\mathrm dx},\ x-w\, .\,\tfrac{\mathrm d\zeta}{\mathrm du},\ \lambda',\lambda'',\lambda'''\text{ etc.})
\]
is identically divisible by $\zeta$. The quotient is a polynomial in the indeterminates $u,x,w,a,b,c$ etc., symmetric with respect to $a,b,c$ etc.; we shall represent it by
\[
\psi(u,x,w,\lambda',\lambda'',\lambda'''\text{ etc.})
\]
We conclude that we also have identically
\[
\begin{aligned}
&f(u+w\, .\,\tfrac{\mathrm dz}{\mathrm dx},\ x-w\, .\,\tfrac{\mathrm dz}{\mathrm du},\ l',l'',l'''\text{ etc.})\\
&=z\psi(u,x,w,l',l'',l'''\text{ etc.})
\end{aligned}
\]
and likewise
\[
\begin{aligned}
&f(u+w\, .\,\tfrac{\mathrm dZ}{\mathrm dx},\ x-w\, .\,\tfrac{\mathrm dZ}{\mathrm du},\ L',L'',L'''\text{ etc.})\\
&=Z\psi(u,x,w,L',L'',L'''\text{ etc.})
\end{aligned}
\]
Thus, if we represent the function $Z$ simply by $F(u,x)$, so that
\[
f(u,x,L',L'',L'''\text{ etc.})=F(u,x)
\]
holds, we have identically
\[
\begin{aligned}
&F(u+w\, .\,\tfrac{\mathrm dZ}{\mathrm dx},\ x-w\, .\,\tfrac{\mathrm dZ}{\mathrm du})\\
&=Z\psi(u,x,w,L',L'',L'''\text{ etc.})
\end{aligned}
\]

\SourceLeaf{061}
\SecondArticle{18}
Suppose, then, that fixed values of $u,x$, namely $u=U,x=X$, give
\[
\frac{\mathrm dZ}{\mathrm dx}=X',\quad\frac{\mathrm dZ}{\mathrm du}=U'
\]
We have identically
\[
\begin{aligned}
&F(U+wX',\ X-wU')\\
&=F(U,X)\, .\,\psi(U,X,w,L',L'',L'''\text{ etc.})
\end{aligned}
\]
Whenever $U'$ does not vanish, we may set
\[
w=\frac{X-x}{U'}
\]
which gives
\[
\begin{aligned}
&F(U+\tfrac{XX'}{U'}-\tfrac{X'x}{U'},\ x)\\
&=F(U,X)\, .\,\psi(U,X,\tfrac{X-x}{U'},L',L'',L'''\text{ etc.})
\end{aligned}
\]
This can also be stated as follows:

If in the function $Z$ we set $u=U+\frac{XX'}{U'}-\frac{X'x}{U'}$, it becomes
\[
F(U,X)\, .\,\psi(U,X,\tfrac{X-x}{U'},L',L'',L'''\text{ etc.})
\]

\SecondArticle{19}
In the case where $P=0$ does not hold, the determinant of the function $Z$ is a function of the indeterminate $x$ that does not vanish in itself. The number of fixed values of $x$ for which this determinant can take the value 0 is therefore clearly finite. Thus infinitely many fixed values of $x$ can be specified which give this determinant a value different from 0. Let $X$ be such a value of $x$, which we may also suppose to be \textit{real}. The determinant of the function $F(u,X)$ is then not $=0$. By Theorem II of Article 6, the functions
\[
F(u,X)\quad\text{and}\quad\frac{\mathrm dF(u,X)}{\mathrm du}
\]
can therefore have no common divisor. Suppose further that there is a fixed value of $u$, namely $U$, either real or imaginary, that is, of the form $g+h\BareRoot{-1}$, which makes $F(u,X)=0$, so that $F(U,X)=0$. Then $u-U$ is a factor of the polynomial $F(u,X)$, and the function $\frac{\mathrm dF(u,X)}{\mathrm du}$ is consequently certainly not divisible by $u-U$. Suppose this function,
\SourceLeaf{062}
$\frac{\mathrm dF(u,X)}{\mathrm du}$, takes the value $U'$ when $u=U$ is set. It certainly cannot be the case that $U'=0$. Now $U'$ is clearly the value of the partial derivative $\frac{\mathrm dZ}{\mathrm du}$ for $u=U,x=X$. If we also denote by $X'$ the value of the partial derivative $\frac{\mathrm dZ}{\mathrm dx}$ for the same values of $u,x$, the preceding article shows that the function $Z$ vanishes identically under the substitution
\[
u=U+\frac{XX'}{U'}-\frac{X'x}{U'}
\]
and is therefore identically divisible by the factor
\[
u+\frac{X'}{U'}x-(U+\frac{XX'}{U'})
\]
Hence, on setting $u=xx$, we see that $F(xx,x)$ is divisible by
\[
xx+\frac{X'}{U'}x-(U+\frac{XX'}{U'})
\]
and thus takes the value 0 when $x$ is taken to be a root of the equation
\[
xx+\frac{X'}{U'}x-(U+\frac{XX'}{U'})=0
\]
that is, when we set
\[
x=\frac{-X'\pm\BareRoot{(4UU'U'+4XX'U'+X'X')}}{2U'}
\]
These values are known to be either real or of the form $g+h\BareRoot{-1}$.

It is now easy to prove that the function $Y$ must also vanish for these same values of $x$. For $f(xx,x,\lambda',\lambda'',\lambda'''\text{ etc.})$ is clearly the product of all $(x-a)(x-b)$, with repetitions excluded, and thus is $=\upsilon^{m-1}$. It follows at once that
\[
\begin{aligned}
f(xx,x,l',l'',l'''\text{ etc.})&=y^{m-1}\\
f(xx,x,L',L'',L'''\text{ etc.})&=Y^{m-1}
\end{aligned}
\]
or $F(xx,x)=Y^{m-1}$, whose fixed value therefore cannot vanish unless the value of $Y$ vanishes at the same time.

\SecondArticle{20}
The preceding investigations have reduced the \textit{solution} of the equation $Y=0$, that is, the finding of a fixed value of $x$, either real or of the form $g+h\BareRoot{-1}$, which satisfies it, to the solution of the equation $F(u,X)=0$,
\SourceLeaf{063}
provided that the determinant of the function $Y$ is not $=0$. It should be observed that, if all the coefficients in $Y$, that is, the numbers $L',L'',L'''$ etc., are real quantities, all the coefficients in $F(u,X)$ are also real, provided that a real quantity is taken for $X$, as is allowed. The degree of the secondary equation $F(u,X)=0$ is given by the number $\frac12m(m-1)$. Whenever $m$ is an even number of the form $2^\mu k$, with $k$ denoting an undetermined odd number, the degree of the secondary equation is therefore given by a number of the form $2^{\mu-1}k$.

If the determinant of the function $Y$ becomes $=0$, Article 10 gives another polynomial $\mathfrak{Y}$ dividing it, whose determinant is not $=0$ and whose degree is given by a number of the form $2^\nu k$, with either $\nu<\mu$ or $\nu=\mu$. Every solution of the equation $\mathfrak{Y}=0$ also satisfies the equation $Y=0$. The solution of the equation $\mathfrak{Y}=0$ is again reduced to the solution of another equation whose degree is given by a number of the form $2^{\nu-1}k$.

We therefore conclude that, in general, the solution of any equation whose degree is an even number of the form $2^\mu k$ can be reduced to the solution of another equation whose degree is of the form $2^{\mu'}k$, with $\mu'<\mu$. Whenever this number is still even, that is, $\mu'$ is not $=0$, the same method is applied again. We shall continue in this way until we reach an equation of odd degree. All its coefficients will be real if all the coefficients of the original equation were real. Such an equation of odd degree is known to have a solution, indeed a real root. Each of the preceding equations will therefore also have a solution, either through real roots or through roots of the form $g+h\BareRoot{-1}$.

It has thus been established that any polynomial $Y$ of the form $x^m-L'x^{m-1}+L''x^{m-2}-$ etc., where $L',L''$ etc. are fixed real quantities, contains a factor $x-A$, where $A$ is a quantity either real or of the form $g+h\BareRoot{-1}$. In the latter case it is easy to see that $Y$ also takes the value 0 under the substitution $x=g-h\BareRoot{-1}$, and is therefore divisible by $x-(g-h\BareRoot{-1})$, hence also by the product $xx-2gx+gg+hh$. Every polynomial $Y$ therefore certainly contains a real factor of the first or second degree. Since the same holds again for the quotient, $Y$ can clearly be resolved into real factors of the first or second degree. This was what the present paper set out to prove.
\begin{center}\rule{3cm}{0.3pt}\end{center}
